% Experiment report. Fixed sections: see experiments/TEMPLATE/report.tex and experiments/README.md.
% Build: npm run experiment:report -- 24-rl-unit-value
% DRAFT by Claude from the design discussion with Jakob on 27 Nov 2026; Jakob edits before the first research run.
\documentclass{../polyexp}
\begin{document}
\input{generated}
\ExpHeader

\section{Question and hypothesis}
% WHEN: before the first run. WHO: Jakob states the idea in his words; Claude turns it into a prediction.
Jakob's idea, in his words (27 Nov 2026): ``What if we train separate value functions for a city and a unit? Then we
could have the unit rewarded by kills and captures [\ldots] Lets start the experiment by just controlling each unit with
its own value function and leaving the rest unchanged.'' The value function must let a unit learn ``the value of
different attack squad structures (which combination of units attack successfully), whom to attack and where it needs
to move.''

Why: curiosity runs on 26--27 Nov showed that value functions over the whole empire, trained in self-play, lose
clearly to the greedy champion (best: 5{,}020 against the champions' 8{,}701, 200 games); they never saw an attack.
A unit decides locally, so a function of the unit's own surroundings can learn fighting, and every unit in every game
is a training sample.

\todo{jakob}{Confirm or change the prediction below.}
Prediction: with the champion playing the economy, the learner's unit play is at least as good as the champion's.
Measured as learner score minus the base bots' mean score in the same game (the champion in the learner seat scores
$-496 \pm 555$ by this measure, i.e.\ about even):
\begin{itemize}
  \item linear: about even with the champion ($-500$ to $+500$); kills and captures per game at least the champion's;
  \item neural net: better than linear by at least 300 points, because squad combinations are interactions a linear
        model cannot represent.
\end{itemize}
Against the idea: the learner clearly below the champion reference after training (worse than $-1{,}000$), or more
kills with a lower score (the reward misses what the score needs).

\section{Mechanism}
% WHEN: before the first run. WHO: Claude.
\textbf{Who acts} (decided with Jakob, 27 Nov: fighting actions to the value function, actions that spend or earn
stars to the bot). The unit value function chooses move, attack, capture, recover, promote, heal others and the tribe
combat abilities. The base bot (\texttt{greedy-full-probe-perf}) plays upgrade, disband, examine (ruins), harvest
starfish and every city, tech and diplomacy action, with its own chooser (\texttt{full\_best\_of}) on the filtered list,
once before and once after the unit phase. The other three seats are the base bot.

\textbf{Choice.} For each legal unit action $a$ of unit $u$: $Q(a) = r(a) + V(\phi_u(s_a))$, where $s_a$ is the state
after $a$ and $\phi_u$ the features of $u$ in it (afterstate values, as in the earlier RL runs). The highest $Q$ is
played; when no unit action beats leaving the units as they are, the turn passes to the end of turn.

\textbf{One function, every unit.} $V$ is shared by all units; each evaluates it from its own local view.

\textbf{Rewards} (in stars):
\begin{itemize}
  \item capture of a city or village: the captured city's income times the turns left in the game;
  \item kill: the victim's cost;
  \item kill credit is split among the learner's units that damaged the victim in the last 3 turns, by damage;
  \item capture credit: half to the capturer, half split by damage dealt within 3 tiles of the city in the last 3
        turns (all to the capturer when nobody fought there);
  \item own city lost: minus its income times the turns left, split evenly among the learner's units;
  \item no explicit death penalty: a dead unit's return ends, so V learns what dying costs from the rewards it loses.
\end{itemize}
Return: Monte Carlo per unit, undiscounted ($\gamma = 1$). Jakob, 27 Nov: ``Do we need gamma? Because there is already a
turn limit'': the game ends at turn 30, so every return is finite, and the capture reward already counts the turns
left; a discount would count that twice and make long marches (across water) look worse than they are. A first pair of
runs with $\gamma = 0.97$ was stopped for this (linear at iteration 399: $-2{,}113$; net at 304: $-2{,}061$; kept as
\texttt{rl-unit-linear-g097}, \texttt{rl-unit-mlp-g097}).

\textbf{Features} of a unit, read after the action:
own unit (attack, defence, range, movement, HP, HP share, cost);
squad within 3 tiles (count, total attack, total HP, number with range $\geq 2$);
nearest enemy unit (attack, defence, HP share, expected damage dealt and received if attacked now, can kill now);
danger (enemy units that can reach the unit next turn, their total attack, defence bonus of its tile);
where (distance to the nearest enemy city, village, enemy unit and threatened own city).

\textbf{Water} (Jakob, 27 Nov: ``what about water?''). On a continents map most attacks cross water, so every distance
is a \emph{travel} distance in turns, found by a search under the unit's own movement rules (land, and water where the
unit can embark or already floats), not a straight line; a target with no route yet gets a flag instead of a distance.
The same search covers obstacles (Jakob: ``and obstacles''): mountains without climbing, extra movement cost of
forests, roads, tiles occupied by units and enemy zones of control, all by the engine's own movement rules, so unit
skills such as climbing count automatically.
A unit on water has its own form (stats read from its current form) and an on-water flag, so its weakness there shows
in its own stats and in the danger features. Embarking and disembarking are ordinary legal actions and are evaluated
like any other move.

\textbf{Two variants}, same features, rewards and seeds: linear ($V = w \cdot \phi$, least squares on all samples,
older iterations decayed by 0.98, ridge 1) and a small neural net (one hidden layer of 32 tanh units, Adam, learning
rate $10^{-3}$, 1{,}500 minibatches of 128 per iteration from a replay buffer of the last 400{,}000 samples).
Exploration: $\varepsilon$-greedy over the unit actions and ``pass'', $\varepsilon$ from 0.3 to 0.05 over the first
200 iterations; ties broken at random.

Code: a new binary crate \texttt{engine/unitrl} (self-play loop on threads, both learners, run file for the RL dashboard).

\section{Setup}
% WHEN: before the first run. WHO: Claude.
Perfection, 4 players, 16$\times$16 continents, 30 turns; all seats human, normal difficulty (the same economic rules
in every seat, so only the unit play differs). Training: the learner in one seat, rotated, against 3$\times$ the base
bot; training seeds from 7{,}000{,}000{,}000. Measurement: 200 games per variant on seeds 6{,}000{,}000{,}000.. (the
seeds of the 27 Nov comparison), the learner seat rotated; the champion in the learner seat on the same seeds is the
reference. The official \texttt{rl-gate} needs a learner kind for this policy (see \#198); until then the 200-game
comparison is the measure, and the gate follows once the policy is playable there.

Runs: 1{,}000 iterations of 128 training games each, plus 64 greedy evaluation games per iteration on the fixed seeds
6{,}000{,}000{,}000..; command \texttt{unit-rl train --run rl-unit-linear|rl-unit-mlp --model linear|mlp --iters 1000
--games 128 --eval 64 --eps-iters 200 --threads 5}. Sanity check: the champion in the learner seat scores 7{,}436
against 7{,}778 on 20 seeds, matching the 27 Nov comparison (7{,}460 against 7{,}956); in a 40-iteration smoke run the
linear learner went from $-2{,}804$ to $-1{,}106$.

Deferred to later iterations (decided 27 Nov): credit for damage absorbed (tanks shielding ranged units), skill flags,
pair features.

\section{Results}
\ExpResults

\todo{claude}{Read the tables.}

\section{Behaviour}
\todo{claude}{Kills, captures and units lost per game, learner vs champion reference.}
\todo{jakob}{What you saw in the live games, if you watched any.}

\section{Interpretation}
\todo{claude}{Why the result came out this way, and whether the prediction held.}
\todo{jakob}{Your reading: agree, disagree, what is missing?}

\section{Threats to validity}
\todo{claude}{Two or three concrete threats.}

\section{Verdict and next step}
\todo{claude}{Kept or not, and the next experiment this suggests.}

\section{Notes for the paper}
\todo{claude}{Draft: the one-sentence claim and which figure or table would carry it.}
\todo{jakob}{Is this worth a place in the paper? Where?}

\end{document}
